% figure (2026-11-17): labels of a general 2 -> 2 process and the three Mandelstam combinations, after PHY 513 Lecture 21 slide 20 and PS p. 156; generated by qft-C/tools/l21/figs/make_figs.sh
\tikzset{
  fdl/.style={line width=0.8pt, black},
  fdx/.style={circle, fill=black, inner sep=0pt, minimum size=3.6pt},
  fdv/.style={circle, fill=figB, draw=figB, inner sep=0pt, minimum size=5pt},
  fdlab/.style={font=\scriptsize, inner sep=2.6pt},
  fdmom/.style={->, >=Stealth, figO, line width=0.6pt},
}
\tikzset{
  fda/.style={fdl, postaction={decorate}, decoration={markings, mark=at position 0.55 with {\arrow{Stealth[length=5pt]}}}},
  fdS/.style={font=\small},
}
\tikzset{
  fline/.style={line width=0.8pt, black, postaction={decorate}, decoration={markings, mark=at position 0.56 with {\arrow{Stealth[length=5.5pt, width=4.5pt]}}}},
  sline/.style={line width=0.8pt, black, dash pattern=on 3.5pt off 2.2pt},
  blob/.style={draw, line width=0.8pt, fill=figB!12},
}
\tikzset{
  pline/.style={line width=0.8pt, black, decorate, decoration={snake, amplitude=1.6pt, segment length=5.2pt}},
}
\begin{tikzpicture}[font=\small]
  \foreach \x/\ttl/\sub in {0/{$s = (p + p')^2$}/{the two incoming momenta}, 5.2/{$t = (k - p)^2$}/{$p$ and the outgoing $k$}, 10.4/{$u = (k' - p)^2$}/{$p$ and the outgoing $k'$}} {
    \begin{scope}[shift={(\x,0)}]
      \draw[fdl] (-1.25,-1.25)--(0,0); \draw[fdl] (1.25,-1.25)--(0,0); \draw[fdl] (0,0)--(-1.25,1.25); \draw[fdl] (0,0)--(1.25,1.25);
      \draw[fdmom] (-1.2,-0.75)--(-0.75,-0.3); \draw[fdmom] (1.2,-0.75)--(0.75,-0.3); \draw[fdmom] (-0.75,0.3)--(-1.2,0.75); \draw[fdmom] (0.75,0.3)--(1.2,0.75);
      \node[circle, blob, minimum size=0.8cm] at (0,0) {};
      \node[fdlab] at (-1.45,-1.42) {$p$}; \node[fdlab] at (1.45,-1.42) {$p'$}; \node[fdlab] at (-1.45,1.42) {$k$}; \node[fdlab] at (1.45,1.42) {$k'$};
      \node[fdlab] at (0,2.25) {\ttl}; \node[font=\scriptsize] at (0,-2.2) {\sub};
    \end{scope}
  }
  % highlight the pair entering each variable
  \draw[figR, line width=1.1pt, rounded corners=6pt] (-1.75,-1.72) rectangle (1.75,-1.1);
  \draw[figR, line width=1.1pt, rounded corners=6pt] (5.2-1.75,-1.72) rectangle (5.2-1.12,1.72);
  \draw[figR, line width=1.1pt, rounded corners=6pt, rotate around={45:(10.4,0)}] (10.4-2.42,-0.32) rectangle (10.4+2.42,0.32);
  \draw[->, gray] (12.9,-0.9)--(12.9,0.9); \node[font=\scriptsize, gray, right] at (12.9,0) {time};
\end{tikzpicture}
